\chapter{Testes de Hipóteses}
\label{cap:hipoteses}

No capítulo anterior, descrevemos as notas observadas. Agora vamos formular perguntas sobre médias populacionais e examinar como os dados se relacionam com essas perguntas. Um teste de hipóteses combina um modelo estatístico, uma hipótese nula e uma estatística de teste. Sua interpretação depende também de como as observações foram obtidas.

Vamos trabalhar com testes t de uma amostra, de duas amostras independentes e de medidas pareadas, avançando para testes de postos, planejamento amostral e multiplicidade. Usaremos novamente \texttt{exemplo\_notas.csv}; todos os resultados numéricos apresentados neste capítulo são cálculos sobre esses dados fictícios.

\section{Hipóteses, valor de p e decisões}

\subsection{Formulando uma pergunta antes da análise}

Suponha que uma nota média de $60$ pontos seja uma referência definida antes de observar os resultados. Podemos perguntar se a média populacional das notas posteriores difere desse valor. O teste bilateral corresponde a:
\[
H_0:\mu=60,\qquad H_1:\mu\neq60.
\]
O parâmetro $\mu$ representa a média na população de interesse; $\bar{x}$ representa a média calculada na amostra. Formular uma hipótese sobre $\mu$ não dispensa definir a população, a estratégia de amostragem e a variável medida.

O valor de $p$ é calculado sob $H_0$ e as demais suposições do modelo: ele expressa a probabilidade de uma estatística de teste tão ou mais extrema que a observada, no sentido especificado pela hipótese alternativa \parencite{NISTValorP}. Antes da análise, escolhemos um nível de significância $\alpha$; se usarmos $\alpha=0{,}05$, um resultado com $p\leq0{,}05$ leva à rejeição de $H_0$ segundo essa regra.

\subsection{O que um teste não permite concluir}

Um valor de $p$ pequeno não informa a probabilidade de $H_0$ ser verdadeira, não mede o tamanho do efeito e não garante importância prática. Um valor de $p$ grande também não demonstra igualdade. A declaração da Associação Americana de Estatística recomenda interpretar os resultados com o desenho do estudo, as estimativas e a incerteza, com transparência sobre as análises realizadas \parencite{WassersteinLazar2016}.

Por exemplo, uma diferença de $0{,}2$ ponto pode ser detectável em uma amostra muito grande e ser pouco relevante para a decisão educacional. Em outra situação, uma diferença estimada de $10$ pontos pode vir acompanhada de um intervalo muito amplo. Relatar somente ``deu significativo'' omite essas informações.

\begin{tcolorbox}[colback=white,colframe=green,title={Dica de Conteúdo}]
Se $p>\alpha$, escreva que não foi encontrada evidência suficiente para rejeitar $H_0$ no teste realizado. Evite escrever que os grupos são iguais. Demonstrar equivalência exige uma pergunta, margens e procedimentos próprios.
\end{tcolorbox}

\subsection{Intervalo de confiança e tamanho de efeito}

Um intervalo de confiança combina estimativa e precisão. Em termos frequentistas, um procedimento de IC95\% cobre o parâmetro verdadeiro em 95\% das repetições do processo de amostragem, se suas suposições forem atendidas. Não significa que 95\% das observações individuais estejam dentro do intervalo da média \parencite{NISTIntervalo}.

Além da diferença em pontos, pode ser útil apresentar uma medida padronizada, como o $d$ de Cohen, oferecida nos testes t do jamovi. Identifique o tipo de comparação e a definição da padronização; o desvio padrão usado em um teste pareado não precisa ser o mesmo de um teste independente. Interprete a magnitude no contexto da variável e da pergunta.

\section{Teste t de uma amostra}

O teste t de uma amostra compara uma média com um valor de referência. Com observações independentes, desvio padrão populacional desconhecido e uma distribuição para a qual a inferência t seja adequada, a estatística é:
\[
t=\frac{\bar{x}-\mu_0}{s/\sqrt{n}},\qquad gl=n-1.
\]
Para os dados de \texttt{depois}, temos $n=18$, $\bar{x}\approx64{,}556$, $s\approx8{,}597$ e $\mu_0=60$. Portanto,
\[
t\approx\frac{64{,}556-60}{8{,}597/\sqrt{18}}\approx2{,}248.
\]

\subsection{Executando no jamovi}

\begin{enumerate}
  \item Desative os filtros e abra \textbf{Análises $\rightarrow$ Testes t $\rightarrow$ Teste t de uma amostra (One Sample T-Test)}.
  \item Mova \texttt{depois} para a lista de variáveis dependentes e informe $60$ em \textbf{Valor do teste (Test value)}.
  \item Mantenha a alternativa bilateral, diferente do valor de referência.
  \item Solicite descritivas, diferença média, intervalo de confiança de 95\% e o gráfico Q-Q para examinar a distribuição.
\end{enumerate}

A documentação oficial descreve as alternativas e as opções de intervalo para esse teste \parencite{JamoviTUmaAmostra}. Você deve encontrar aproximadamente $t(17)=2{,}248$ e $p=0{,}038$. A diferença estimada é $4{,}556$ pontos, com IC95\% de $[0{,}280;8{,}831]$. Como o intervalo da saída se refere à \textbf{diferença em relação a 60}, somamos $60$ a seus limites para obter o intervalo da média: $[60{,}280;68{,}831]$.

Um relato possível é: ``No exemplo fictício, a média posterior foi de $64{,}56$ pontos. A diferença em relação à referência de $60$ foi estimada em $4{,}56$ pontos, IC95\% $[0{,}28;8{,}83]$, $t(17)=2{,}25$, $p=0{,}038$.'' A referência não representa, por si só, uma medida anterior das mesmas pessoas.

\section{Teste de hipóteses unicaudal e bilateral}

Uma alternativa unilateral especifica uma direção. Se a pergunta planejada for se a média é maior que $60$, podemos escrever $H_0:\mu\leq60$ e $H_1:\mu>60$; o cálculo é feito na fronteira $\mu=60$. Para uma alternativa de redução, as desigualdades são invertidas. No teste bilateral, desvios nas duas direções contam contra a hipótese nula.

A direção deve resultar da pergunta de pesquisa, antes de examinar os dados. Não troque para um teste unilateral apenas porque o valor de $p$ bilateral ficou acima de $0{,}05$. Se uma diferença na direção oposta também for relevante para a decisão, explique por que uma análise unilateral seria apropriada antes de usá-la.

No jamovi, escolha a hipótese alternativa no painel do teste e confira a nota da tabela. Em testes de dois grupos, a interpretação de ``maior que'' depende da ordem dos níveis; em testes pareados, depende da ordem das medidas. Sempre escreva a diferença que está sendo testada.

\section{Duas amostras independentes: Student e Welch}

Quando cada pessoa pertence a apenas um dos dois grupos comparados, a estrutura é independente. Nosso exemplo permite comparar Leitura e Vídeo, usando a nota \texttt{depois}. As notas anterior e posterior da mesma pessoa, por outro lado, não formam dois grupos independentes.

O teste t de Student com variância combinada supõe variâncias populacionais iguais. O teste de Welch dispensa essa igualdade e usa graus de liberdade aproximados. Em ambos, a independência vem do desenho do estudo. A análise deve considerar a forma das distribuições, o tamanho das amostras e possíveis valores influentes \parencite{NISTTesteT,JamoviTIndependente}.

\subsection{Configurando a comparação}

\begin{enumerate}
  \item Crie um filtro com \texttt{grupo != 'Oficina'}, mantendo apenas Leitura e Vídeo. Confira que restaram 12 observações.
  \item Abra \textbf{Testes t $\rightarrow$ Teste t de amostras independentes (Independent Samples T-Test)}.
  \item Atribua \texttt{depois} à variável dependente e \texttt{grupo} à variável de agrupamento.
  \item Selecione \textbf{Welch}, descritivas, diferença média, IC95\% e gráficos de diagnóstico. Confira a ordem dos níveis antes de interpretar o sinal.
\end{enumerate}

Para a diferença \textbf{Leitura menos Vídeo}, a estimativa é $-5{,}50$ pontos. O teste bilateral de Welch produz aproximadamente $t=-1{,}287$, $gl=9{,}902$ e $p=0{,}227$, com IC95\% $[-15{,}037;4{,}037]$. Se a ordem for invertida, o sinal da diferença e de $t$ também será invertido.

O intervalo é compatível com diferenças em ambos os sentidos. Com seis pessoas por grupo, esse exemplo não demonstra equivalência nem permite afirmar que as atividades são intercambiáveis. \textbf{Desative o filtro} antes de continuar para o teste pareado.

\section{Amostras pareadas: antes e depois}

No teste pareado, trabalhamos com uma diferença para cada pessoa. Definindo $d_i=\text{depois}_i-\text{antes}_i$, testamos se a média populacional dessas diferenças é zero:
\[
t=\frac{\bar{d}}{s_d/\sqrt{n}},\qquad gl=n-1.
\]
A unidade da análise é o par. A suposição de normalidade, quando usada para a inferência t exata, diz respeito às \textbf{diferenças}, e não à obrigação de as duas colunas apresentarem separadamente uma distribuição normal \parencite{JamoviTPareado}.

\subsection{Configurando e interpretando}

\begin{enumerate}
  \item Com os 18 casos disponíveis, abra \textbf{Testes t $\rightarrow$ Teste t de amostras pareadas (Paired Samples T-Test)}.
  \item Forme um par com \texttt{depois} como primeira medida e \texttt{antes} como segunda.
  \item Solicite a diferença média, IC95\%, descritivas e diagnóstico das diferenças.
  \item Confira os identificadores e o número de pares completos. Nunca ordene uma das colunas separadamente.
\end{enumerate}

Para \texttt{depois} menos \texttt{antes}, a diferença média é $12{,}889$ pontos, com desvio padrão das diferenças de $5{,}697$. O resultado é aproximadamente $t(17)=9{,}598$, $p<0{,}001$, IC95\% $[10{,}056;15{,}722]$. São 18 pares, não 36 observações independentes.

Esse resultado descreve uma mudança no exemplo fictício. Em um estudo real, uma comparação antes/depois sem controle adequado pode refletir outros fatores, como familiarização com a prova ou passagem do tempo. O teste, isoladamente, não estabelece a causa da mudança.

\section{Pressupostos e alternativas}

Antes de interpretar um teste, confira como as pessoas foram selecionadas, se existem medidas repetidas e como os ausentes foram tratados. Um teste de normalidade não verifica independência. Seu resultado também não deve funcionar como uma chave automática que escolhe toda a análise: em amostras pequenas ele pode ter pouca capacidade de detectar desvios, e em amostras grandes pode detectar diferenças pouco relevantes para a pergunta.

O jamovi oferece Mann--Whitney para amostras independentes e Wilcoxon de postos sinalizados para situações pareadas ou de uma amostra. Esses testes têm hipóteses e condições próprias. Mann--Whitney não é, de modo geral, simplesmente um teste de diferença entre medianas; interpretações dessa forma exigem condições adicionais sobre as distribuições. Wilcoxon de postos sinalizados também exige atenção à simetria das diferenças para uma interpretação como teste de localização.

Quando houver dúvida, volte à pergunta e examine gráficos e suposições. Registre a análise escolhida, as exclusões, o número de observações, a estimativa, seu intervalo e o valor de $p$. A conclusão deve responder à pergunta planejada, sem escolher apenas o resultado mais favorável entre várias tentativas.

\section{Escolhendo a análise a partir do desenho}
\label{sec:escolha-teste}

Vamos organizar a decisão antes de abrir o menu. Primeiro, escreva qual é a variável resposta e o que uma linha representa. Depois, identifique se as observações vêm de pessoas distintas ou se a mesma pessoa aparece em mais de uma condição. Por fim, explicite a quantidade que queremos estimar: uma média, uma diferença média, uma proporção ou uma associação.

\begin{table}[H]
\centering
\caption{Perguntas iniciais para orientar a análise}
\begin{tabular}{p{0.30\textwidth}p{0.27\textwidth}p{0.32\textwidth}}
\hline
Pergunta & Estrutura & Ponto de partida \\
\hline
Média versus referência & Uma medida por pessoa & t de uma amostra \\
Duas médias & Dois grupos independentes & t independente; considerar Welch \\
Mudança média & Duas medidas da mesma pessoa & t pareado nas diferenças \\
Mais de duas médias & Grupos independentes & ANOVA de um fator \\
Evolução em três ocasiões & Medidas das mesmas pessoas & ANOVA de medidas repetidas \\
Percentual versus referência & Resposta binária & Teste binomial \\
\hline
\end{tabular}
\label{tab:perguntas-testes}
\end{table}

A tabela não encerra a escolha. Por exemplo, pessoas agrupadas em escolas podem continuar dependentes mesmo que cada pessoa tenha uma única linha. Uma escala ordinal com poucas categorias também não deve ser tratada automaticamente como uma medida contínua. Nos capítulos seguintes, retomaremos o desenho para escolher modelos com fatores, covariáveis e respostas categóricas.

\section{Mann--Whitney e Wilcoxon na prática}

\subsection{Dois grupos independentes}

Para praticar Mann--Whitney, mantenha somente Leitura e Vídeo em \texttt{exemplo\_notas.csv}, como na comparação de Welch. Abra o teste t de amostras independentes, atribua \texttt{depois} à resposta e \texttt{grupo} ao agrupamento e selecione a opção \textbf{Mann--Whitney U}. Acrescente as descritivas por grupo e examine a distribuição dos valores \parencite{JamoviTIndependente}.

O procedimento usa a ordenação conjunta das observações. Valores iguais recebem postos médios. Sob a hipótese nula usual de distribuições iguais, os rótulos dos grupos são intercambiáveis entre observações independentes; a interpretação como deslocamento de localização pede condições adicionais. Se as formas e dispersões forem diferentes, um resultado pode refletir essas diferenças, sem uma interpretação exclusiva em termos de medianas.

Na leitura da saída, registre o tamanho de cada grupo, a estatística, a alternativa e o método de cálculo do valor de $p$. As convenções para a estatística de postos variam: uma soma de postos e um $U$ podem codificar a mesma ordenação com números diferentes. A documentação do R discute essa diferença e distingue a estimativa de localização da simples subtração de medianas \parencite{RWilcoxon}.

\subsection{Duas medidas pareadas}

Desative o filtro e abra o teste de amostras pareadas. Forme o par \texttt{depois} menos \texttt{antes} e selecione \textbf{Wilcoxon de postos sinalizados (Wilcoxon signed-rank)} \parencite{JamoviTPareado}. O procedimento ordena as magnitudes das diferenças não nulas e considera seus sinais; ele preserva a ligação entre as duas medidas de cada pessoa.

No nosso arquivo, todas as 18 diferenças são positivas. Sem diferenças nulas, os postos de 1 a 18 somam $18\times19/2=171$: toda a soma está do lado positivo. Esse cálculo ajuda a entender a estatística; não substitui a identificação da convenção e do método de inferência adotados na saída. Empates, zeros e aproximações podem afetar o cálculo. Para interpretar o procedimento como um teste de localização das diferenças, precisamos também considerar sua simetria \parencite{RWilcoxon}.

Um relato deve indicar que as medidas são pareadas, apresentar uma descrição das diferenças e identificar o teste. Não escolha entre t pareado e Wilcoxon apenas pelo menor valor de $p$. Defina a pergunta e a estratégia principal e, se fizer uma análise de sensibilidade, explique sua finalidade.

\section{Erros, poder e planejamento amostral}

O erro do tipo I consiste em rejeitar uma hipótese nula verdadeira. O erro do tipo II consiste em deixar de rejeitá-la diante de uma alternativa específica verdadeira. Chamamos de poder a probabilidade de rejeição nessa alternativa, $1-\beta$. Não existe um único poder de um teste sem especificar o tamanho do efeito e as demais condições.

Considere um estudo futuro com dois grupos independentes, alocação igual e uma diferença de interesse de 5 pontos. Se o desvio padrão comum planejado for 10 pontos, a diferença padronizada será $d=5/10=0{,}5$. Para um teste t bilateral com $\alpha=0{,}05$ e poder de 80\%, o cálculo convencional indica aproximadamente \textbf{64 pessoas por grupo}, antes de considerar perdas. A função oficial \texttt{power.t.test} do R permite reproduzir o planejamento informando diferença, desvio padrão, poder, nível e tipo de teste \parencite{RPoderT}.

Esse número depende das suposições. Não o reutilize para um estudo pareado, por conglomerados ou com outro desfecho. Documente de onde vieram a diferença relevante e o desvio padrão. Quando não há boa informação prévia, compare cenários plausíveis e inclua o tratamento de perdas no planejamento. Os seis casos por grupo da nossa base foram escolhidos para facilitar a prática; não representam uma recomendação de amostra.

Depois de coletar os dados, examine principalmente as estimativas e seus intervalos. Um resultado impreciso continua impreciso, mesmo que receba uma explicação baseada em ``baixo poder''. Planejar a precisão antes da coleta é mais útil do que tentar justificar a conclusão após observar o valor de $p$.

\section{Múltiplas perguntas e ajustes}

Se fizermos 20 testes independentes, todos com hipóteses nulas verdadeiras e nível de 5\%, a probabilidade de pelo menos uma rejeição será $1-(1-0{,}05)^{20}\approx64\%$. A conta pressupõe independência; o problema da multiplicidade, porém, também existe quando os testes são dependentes. Precisamos definir quais testes pertencem à mesma família de perguntas.

Bonferroni compara cada valor de $p$ com $\alpha/m$, ou multiplica cada valor por $m$, limitando o resultado a 1. Holm ordena os valores e aplica um ajuste sequencial que controla a probabilidade de pelo menos um erro do tipo I na família. Por exemplo, para três valores brutos $0{,}010$, $0{,}030$ e $0{,}040$, os valores ajustados de Holm são $0{,}030$, $0{,}060$ e $0{,}060$ \parencite{RAjusteP}.

Procedimentos voltados à taxa de falsas descobertas, como Benjamini--Hochberg, têm outro objetivo de controle e condições próprias. Eles não significam que cada resultado rejeitado tenha, individualmente, uma probabilidade de 5\% de ser falso. A documentação de \texttt{p.adjust} distingue essas famílias de procedimentos \parencite{RAjusteP}. No jamovi, use as opções de ajuste oferecidas pela análise correspondente e informe o método; não suponha que uma escolha em um painel ajustará automaticamente todos os testes da sessão.

\section{Vamos praticar: perguntas e comentários}

Antes de consultar os comentários, escreva uma resposta curta para cada situação.

\begin{enumerate}
  \item Uma turma faz a mesma avaliação em dois momentos. Podemos usar um teste independente porque existem duas colunas?
  \item O teste de Welch produziu $p=0{,}227$ para Leitura versus Vídeo. Isso demonstra equivalência?
  \item Um teste de Mann--Whitney apresentou evidência de diferença. Podemos escrever automaticamente que as medianas populacionais diferem?
  \item Em três testes planejados, os valores de $p$ foram $0{,}010$, $0{,}030$ e $0{,}040$. Quais rejeitamos a 5\% usando Holm?
\end{enumerate}

\subsection*{Comentários das respostas}

\begin{enumerate}
  \item Não. São medidas das mesmas pessoas; precisamos preservar os pares. Para uma pergunta sobre mudança média, o t pareado trabalha com as diferenças individuais.
  \item Não. O IC95\% da diferença, $[-15{,}037;4{,}037]$, inclui valores em ambos os sentidos. Uma pergunta de equivalência exige margens justificadas e análise própria.
  \item Não sem condições adicionais. Precisamos examinar o desenho e a forma das distribuições e identificar a quantidade que o teste permite interpretar.
  \item Apenas a primeira hipótese, pois os valores ajustados são $0{,}030$, $0{,}060$ e $0{,}060$. A família deve incluir os testes pertinentes, inclusive os que não tiveram resultados favoráveis.
\end{enumerate}
